\documentclass[10pt,a4paper]{amsart}
\usepackage[margin=19mm]{geometry}
\usepackage{amsmath,amssymb,mathtools}
\usepackage[scr=rsfs,cal=euler]{mathalfa}
\usepackage{xfrac}
\usepackage[unicode,hidelinks,bookmarks=false]{hyperref}
\newcommand{\ie}{i.e.\ }
\newcommand{\cf}{cf.\ }
\newcommand{\resp}{respectively\ }
\newcommand{\Subsection}{\subsection}
\providecommand{\nobreakdash}{\nobreak\hbox{-}}
\setlength{\emergencystretch}{2em}
\multlinegap0pt
\usepackage{amssymb}
\usepackage{algorithm}\usepackage{algpseudocode}
\usepackage{bm}
\usepackage{mathtools}
\newtheorem{proposition}{Proposition}
\title{Employing Deep Neural Operators for PDE control by decoupling training and optimization}
\author{Oliver G. S. Lundqvist, Fabricio Oliveira}
\date{Source: arXiv:2506.04742v3 (CC BY 4.0)}
\begin{document}
\maketitle

\noindent\emph{Source and licence.} Employing Deep Neural Operators for PDE control by decoupling training and optimization. Version arXiv:2506.04742v3 (CC BY 4.0), distributed by arXiv under the Creative Commons Attribution 4.0 International licence, http://creativecommons.org/licenses/by/4.0/, as recorded in the arXiv licence metadata for this exact version and shown on the abstract page https://arxiv.org/abs/2506.04742v3. Full licence notice: \texttt{licenses/arxiv-2506.04742-CC-BY-4.0.txt}.\par\smallskip

\noindent\textit{Editorial scope.} Counted unit: the article's proposition prop:cheating\_directions (source0.tex lines 524--533). The article's complete original proof of it is delivered with it (source0.tex lines 535--601, one \texttt{proof} environment of 67 lines). No mathematics is written, repaired or re-derived: the statement and the proof are the article's own lines, in the article's own notation, and no retained line is substituted or re-flowed.  No further source-local context is needed: every cross-reference of every retained line resolves inside the two delivered spans.  Nothing was omitted from any retained span, so the package carries no omission record.  No retained line contains a citation, so the wrapper carries no bibliography and no bibliography entry.  No retained line contains a figure environment, an image-inclusion command or drawing code, so no image is delivered and the package declares no asset dependency.  Editorial formatting only, in addition: an AMS \texttt{amsart} preamble that restates the article's own declarations and macros used by the retained lines, namely the environments \texttt{proposition} on separate counters, together with the standard packages those lines need.  The retained environments print in this document as Proposition 1 in order of appearance, whereas the article prints the same items with the numbering of its own full document.  The complete native source of the article as distributed in the arXiv e-print for this exact version is bundled unchanged as \texttt{sources/179358/source0.tex}.
\begin{proposition}[Existence of cheating directions]\label{prop:cheating_directions}
Let $R(\mathbf{u}_S)=\|\mathcal{R}_\theta(\mathbf{u}_S)\|^2$ and define $\nabla_\mathbf{u}\bar{J}_{NO}(\mathbf{u}_S):=\mathbf{g}$ and $\nabla_\mathbf{u}R(\mathbf{u}_S):=\mathbf{r}$ as the gradients of the cost function and the residual penalty term, respectively. If $\mathbf{g}\neq0$, $\mathbf{r}\neq0$, and $\mathbf{g}$ and $\mathbf{r}$ are not collinear, i.e. $\mathbf{g}\neq c\mathbf{r}$ for some constant $c\neq0$, then there exists a cheating direction $\mathbf{d}$ such that
\begin{equation}\label{eq:cheating_direction}
\mathbf{d}^\top \mathbf{g}<0 \quad \text{and}\quad \mathbf{d}^\top \mathbf{r}>0.
\end{equation}
Consequently, for sufficiently small $\eta>0$, it follows that
\begin{equation}\label{eq:cheating_direction_step}
\bar{J}_{NO}(\mathbf{u}_S+\eta \mathbf{d})<\bar{J}_{NO}(\mathbf{u}_S) \quad \text{and}\quad R(\mathbf{u}_S+\eta \mathbf{d})>R(\mathbf{u}_S).
\end{equation}
\end{proposition}

\begin{proof}
We consider three cases based on the sign of 
$\mathbf{g}^\top\mathbf{r}$.


\medskip\noindent
\textbf{Case 1: $\mathbf{g}^\top \mathbf{r}>0$.}
Consider the family of directions
\begin{equation}\label{eq:family_pos}
\mathbf{d}=-\mathbf{g}+\alpha \mathbf{r}, \quad \alpha\in \mathbb{R}.
\end{equation}
Then
\begin{equation}
\mathbf{g}^\top \mathbf{d}=\mathbf{g}^\top(-\mathbf{g}+\alpha \mathbf{r})=-\|\mathbf{g}\|^2+\alpha \mathbf{g}^\top \mathbf{r},
\end{equation}
and similarly
\begin{equation}
\mathbf{r}^\top \mathbf{d}=\mathbf{r}^\top(-\mathbf{g}+\alpha \mathbf{r})=-\mathbf{r}^\top \mathbf{g}+\alpha\|\mathbf{r}\|^2
=-(\mathbf{g}^\top \mathbf{r})+\alpha\|\mathbf{r}\|^2.
\end{equation}
Thus $\mathbf{g}^\top \mathbf{d}<0$ holds whenever $\alpha<\|\mathbf{g}\|^2/(\mathbf{g}^\top \mathbf{r})$, and $\mathbf{r}^\top \mathbf{d}>0$ holds whenever
$\alpha>(\mathbf{g}^\top \mathbf{r})/\|\mathbf{r}\|^2$. Such an $\alpha$ exists if and only if
\begin{equation}
\frac{\mathbf{g}^\top \mathbf{r}}{\|\mathbf{r}\|^2}<\frac{\|\mathbf{g}\|^2}{\mathbf{g}^\top \mathbf{r}}
\iff (\mathbf{g}^\top \mathbf{r})^2 <\|\mathbf{g}\|^2\|\mathbf{r}\|^2,
\end{equation}
which holds by the strict Cauchy--Schwarz inequality since $\mathbf{g}$ and $\mathbf{r}$ are not collinear.
Therefore there exists $\alpha$ such that $\mathbf{g}^\top \mathbf{d}<0$ and $\mathbf{r}^\top \mathbf{d}>0$.

\medskip\noindent
\textbf{Case 2: $\mathbf{g}^\top \mathbf{r} < 0$.}
Take $\mathbf{d} = -\mathbf{g}$. Then
\begin{equation}
\mathbf{g}^\top\mathbf{d} = -\|\mathbf{g}\|^2 < 0,
\end{equation}
and
\begin{equation}
\mathbf{r}^\top\mathbf{d} = -\mathbf{g}^\top\mathbf{r} > 0.
\end{equation}
That is, the steepest descent direction for the tracking cost is itself
a cheating direction.

\medskip\noindent
\textbf{Case 3: $\mathbf{g}^\top \mathbf{r} = 0$.}
Take $\mathbf{d} = -\mathbf{g} + \alpha\mathbf{r}$ with any $\alpha > 0$. Then
\begin{equation}
\mathbf{g}^\top\mathbf{d} = -\|\mathbf{g}\|^2 < 0,
\end{equation}
and
\begin{equation}
\mathbf{r}^\top\mathbf{d} = \alpha\|\mathbf{r}\|^2 > 0.
\end{equation}




\medskip\noindent
Finally, since $\bar{J}_{\mathrm{NO}}$ and $R$ are differentiable, for sufficiently small $\eta>0$ we have
\[
\bar{J}_{\mathrm{NO}}(\mathbf{u}_S+\eta \mathbf{d})=\bar{J}_{\mathrm{NO}}(\mathbf{u}_S)+\eta\, \mathbf{g}^\top \mathbf{d}+o(\eta),
\qquad
R(\mathbf{u}_S+\eta \mathbf{d})=R(\mathbf{u}_S)+\eta\, \mathbf{r}^\top \mathbf{d}+o(\eta),
\]
so $\mathbf{g}^\top \mathbf{d}<0$ and $\mathbf{r}^\top \mathbf{d}>0$ imply
$\bar{J}_{\mathrm{NO}}(\mathbf{u}_S+\eta \mathbf{d})<\bar{J}_{\mathrm{NO}}(\mathbf{u}_S)$ and
$R(\mathbf{u}_S+\eta \mathbf{d})>R(\mathbf{u}_S)$ for all sufficiently small $\eta>0$.
\end{proof}

\end{document}
